\section{Problem Setup}

\subsection{Task}

The system receives a user question about endurance training, training structure, safety, risk, or evidence. It must return one of three output states:

\begin{itemize}[leftmargin=*]
  \item \texttt{answered}: the system provides an evidence-supported answer and passes audit;
  \item \texttt{partial\_answer}: the system provides a bounded answer after \texttt{citation repair};
  \item \texttt{refused}: the system declines to answer because the evidence is insufficient, the request is unsafe, or the request asks for unsupported claims, fabricated evidence, or unsafe continuation.
\end{itemize}

\subsection{Evidence Constraint}

The workflow must keep claims traceable to retrieved or curated gold evidence. If evidence is insufficient, refusal is a valid output. The system should not invent paper titles, page numbers, DOI-like identifiers, or precise performance guarantees.

\subsection{Safety Constraint}

Safety-sensitive requests include red-flag symptoms, worsening pain, fever, suspected infection, heat illness, cardiovascular warning signs, excessive fatigue, or pressure to continue training despite risk. In such cases, the safe behavior may be to refuse or de-escalate rather than continue with a training prescription.

\subsection{Stress Types}

The safety-stress suite targets five request patterns:

\begin{itemize}[leftmargin=*]
  \item prompt injection;
  \item unsafe request;
  \item citation hallucination;
  \item overclaim request;
  \item evidence conflict / cherry-picking.
\end{itemize}
